\documentclass[11pt]{article}
\usepackage[margin=1in]{geometry}
\usepackage{amsmath, amssymb, amsthm}
\usepackage{enumitem}
\usepackage{xcolor}
\usepackage{tcolorbox}
\tcbuselibrary{skins, breakable}

% Define color boxes for different framework sections
\definecolor{problemconfig}{RGB}{230, 242, 255}
\definecolor{strategic}{RGB}{255, 230, 242}
\definecolor{tactical}{RGB}{242, 255, 230}
\definecolor{techniques}{RGB}{255, 242, 230}
\definecolor{verification}{RGB}{255, 230, 230}
\definecolor{solution}{RGB}{230, 255, 255}
\definecolor{competition}{RGB}{242, 242, 255}
\definecolor{gold}{RGB}{218, 165, 32}

\newtcolorbox{problemconfigbox}{
    colback=problemconfig, colframe=blue!40!black, title=PROBLEM CONFIGURATION ANALYSIS,
    fonttitle=\bfseries, breakable
}

\newtcolorbox{strategicbox}{
    colback=strategic, colframe=pink!40!black, title=STRATEGIC FRAMEWORK APPLICATION,
    fonttitle=\bfseries, breakable
}

\newtcolorbox{tacticalbox}{
    colback=tactical, colframe=green!40!black, title=TACTICAL METHODOLOGY SELECTION,
    fonttitle=\bfseries, breakable
}

\newtcolorbox{techniquesbox}{
    colback=techniques, colframe=orange!40!black, title=CONVERSION TECHNIQUES APPLICATION,
    fonttitle=\bfseries, breakable
}

\newtcolorbox{verificationbox}{
    colback=verification, colframe=red!40!black, title=VERIFICATION STRATEGY DESIGN,
    fonttitle=\bfseries, breakable
}

\newtcolorbox{solutionbox}{
    colback=solution, colframe=cyan!40!black, title=SOLUTION ROADMAP,
    fonttitle=\bfseries, breakable
}

\newtcolorbox{competitionbox}{
    colback=competition, colframe=purple!40!black, title=COMPETITION STRATEGY INTEGRATION,
    fonttitle=\bfseries, breakable
}

\newtcolorbox{keyinsightbox}{
    colback=yellow!20, colframe=gold, title=KEY INSIGHT,
    fonttitle=\bfseries, arc=3pt, boxrule=2pt, leftrule=5pt
}

\newtcolorbox{warningbox}{
    colback=red!10, colframe=red!50, title=WARNING: Common Pitfall,
    fonttitle=\bfseries, arc=3pt, boxrule=2pt, leftrule=5pt
}

\newtcolorbox{tipbox}{
    colback=green!10, colframe=green!50, title=PRO TIP,
    fonttitle=\bfseries, arc=3pt, boxrule=2pt, leftrule=5pt
}

\title{\textbf{2017 AMC 12B Problem 19: Strategic Analysis}}
\author{Using AMC12/COMC Problem-Solving Framework}
\date{}

\begin{document}
\maketitle

\section*{Problem Statement}
Let $N=123456789101112\dots4344$ be the $79$-digit number that is formed by writing the integers from $1$ to $44$ in order, one after the other. What is the remainder when $N$ is divided by $45$?

\textbf{Answer Choices:}\\
$\textbf{(A)}\ 1\qquad\textbf{(B)}\ 4\qquad\textbf{(C)}\ 9\qquad\textbf{(D)}\ 18\qquad\textbf{(E)}\ 44$

\begin{problemconfigbox}
\subsection*{A. Problem Classification}
\begin{itemize}[leftmargin=*]
    \item \textbf{Problem Type}: To Find -- determine a specific remainder value
    \item \textbf{Difficulty Band}: Late-band (Problem 19 of 25) -- requires insight into modular arithmetic and Chinese Remainder Theorem
    \item \textbf{Competition Context}: AMC12 (75 min, 25 problems) -- approximately 3-4 minutes allocated for late-band problems
    \item \textbf{Mathematical Domain}: Number Theory (modular arithmetic, divisibility, Chinese Remainder Theorem)
\end{itemize}

\subsection*{B. Problem Structure Identification}
\begin{itemize}[leftmargin=*]
    \item \textbf{Given Information}: 
    \begin{itemize}
        \item $N$ is formed by concatenating integers from 1 to 44
        \item $N = 123456789101112\dots4344$ (79-digit number)
        \item Need remainder when divided by 45
    \end{itemize}
    
    \item \textbf{Constraints}: 
    \begin{itemize}
        \item Must handle a number too large for direct computation
        \item Divisor is $45 = 9 \times 5 = 3^2 \times 5$
        \item Answer must be in range $[0, 44]$
    \end{itemize}
    
    \item \textbf{Objective}: Find $N \bmod 45$
    
    \item \textbf{Hidden Assumptions}: 
    \begin{itemize}
        \item Since $\gcd(9,5) = 1$, can use Chinese Remainder Theorem
        \item Sum of digits gives remainder mod 9
        \item Last digit gives remainder mod 5
        \item Position values matter for mod 45 but not for mod 9
    \end{itemize}
    
    \item \textbf{Answer Choice Leverage}: 
    \begin{itemize}
        \item All choices are $< 45$
        \item Choices 1, 4, 9, 18, 44 suggest testing divisibility patterns
        \item 9 and 18 are multiples of 9; 44 ends in 4
        \item Can eliminate choices that fail either mod 9 or mod 5 test
    \end{itemize}
\end{itemize}

\subsection*{C. Strategic Orientation}
\begin{itemize}[leftmargin=*]
    \item \textbf{Hypothesis}: $N$ is too large to compute directly, must use modular arithmetic
    
    \item \textbf{Conclusion}: Determine $N \bmod 45$ using CRT by finding $N \bmod 9$ and $N \bmod 5$
    
    \item \textbf{Notation}: 
    \begin{itemize}
        \item $N \equiv r \pmod{45}$ where $0 \leq r < 45$
        \item $N \equiv a \pmod{9}$ and $N \equiv b \pmod{5}$
        \item Digit sum $S = 1+2+\cdots+44 = \frac{44 \cdot 45}{2} = 990$
    \end{itemize}
    
    \item \textbf{Intuitive Approach}: 
    \begin{itemize}
        \item Since $45 = 9 \times 5$, decompose the problem
        \item For mod 9: sum of all digits
        \item For mod 5: last digit only
        \item Recombine using CRT
    \end{itemize}
    
    \item \textbf{Cross-Domain Potential}: Could visualize as digit representation in base 10, but modular arithmetic is most direct
\end{itemize}
\end{problemconfigbox}

\begin{strategicbox}
\subsection*{A. Psychological Strategies}
\begin{itemize}[leftmargin=*]
    \item \textbf{Mental Toughness}: Don't be intimidated by the 79-digit number; recognize that modular arithmetic makes it tractable
    
    \item \textbf{Creativity Cultivation}: Consider multiple approaches (direct mod 45, CRT decomposition, pattern recognition); CRT is the most elegant
    
    \item \textbf{Iterative Refinement}: If direct computation fails, decompose into mod 9 and mod 5; verify each step
\end{itemize}

\subsection*{B. Investigation Strategies}
\begin{itemize}[leftmargin=*]
    \item \textbf{Startup Orientation}: Recognize this as a modular arithmetic problem with composite modulus -- immediate cue for CRT
    
    \item \textbf{Get Your Hands Dirty}: Write out first few terms to understand structure: $1,2,3,\ldots,9,10,11,\ldots$
    
    \item \textbf{Penultimate Step}: Once we have $N \equiv 0 \pmod{9}$ and $N \equiv 4 \pmod{5}$, CRT gives the answer
    
    \item \textbf{Wishful Thinking}: We wish the modulus were prime (making this easier), but since $45 = 9 \times 5$, CRT is the natural tool
    
    \item \textbf{Make It Easier}: Instead of computing $N$ directly, compute $N \bmod 9$ and $N \bmod 5$ separately
    
    \item \textbf{Conjecture via Small Cases}: Test with smaller concatenations: $N_3 = 123$, find $N_3 \bmod 45$; observe patterns
    
    \item \textbf{Visualization}: Imagine the number as a sequence of digits; for mod 9, only their sum matters
\end{itemize}

\subsection*{C. Late-Band Strategic Playbook}
\begin{itemize}[leftmargin=*]
    \item \textbf{Answer-Choice Leverage}: Test each answer: must be $\equiv 0 \pmod{9}$ (from digit sum) and $\equiv 4 \pmod{5}$ (from last digit)
    
    \item \textbf{Make It Easier}: Simplify to finding remainders mod 9 and mod 5 separately
    
    \item \textbf{Penultimate Step}: The last step is applying CRT to combine $N \equiv 0 \pmod{9}$ and $N \equiv 4 \pmod{5}$
    
    \item \textbf{Wishful Thinking}: If we knew $N \bmod 9$ and $N \bmod 5$, CRT immediately gives $N \bmod 45$
    
    \item \textbf{Conjecture via Small Cases}: Compute digit sums for smaller ranges to verify pattern
\end{itemize}

\subsection*{D. Argument Construction Strategy}
\begin{itemize}[leftmargin=*]
    \item \textbf{Direct Proof}: Use CRT directly: since $\gcd(9,5)=1$, find $N \bmod 9$ and $N \bmod 5$, then combine
    
    \item \textbf{Proof by Contradiction}: Not applicable here
    
    \item \textbf{Mathematical Induction}: Could use for digit sum formula, but unnecessary
    
    \item \textbf{Stylistic Conventions}: WLOG not needed; direct computation is cleanest
\end{itemize}

\subsection*{E. Cross-Domain Translation}
\begin{itemize}[leftmargin=*]
    \item \textbf{Draw a Picture}: Could visualize digit sequence, but algebraic approach is more efficient
    
    \item \textbf{Recast the Problem}: Instead of ``79-digit concatenation,'' think ``sum of specific powers of 10 with specific coefficients''
    
    \item \textbf{Change Your Point of View}: View as sum of digits (for mod 9) and last digit (for mod 5)
\end{itemize}
\end{strategicbox}

\begin{tacticalbox}
\subsection*{A. Core Mathematical Tactics}
\begin{itemize}[leftmargin=*]
    \item \textbf{Symmetry}: Not directly applicable
    
    \item \textbf{Extreme Principle}: Not applicable
    
    \item \textbf{Invariants}: The remainder classes mod 9 and mod 5 are invariant under CRT
    
    \item \textbf{Monovariants}: Not applicable
    
    \item \textbf{Parity}: Divisibility by 9 (even sum of digits) and by 5 (last digit)
    
    \item \textbf{Graph Theory}: Not applicable
    
    \item \textbf{Generating Functions}: Not applicable
    
    \item \textbf{Linear Algebra}: Not applicable
\end{itemize}

\subsection*{B. Crossover Techniques}
\begin{itemize}[leftmargin=*]
    \item \textbf{Algebraic-Geometric}: Not applicable
    
    \item \textbf{Geometric-Probabilistic}: Not applicable
    
    \item \textbf{Combinatorial-Algebraic}: Not applicable
    
    \item \textbf{Number-Theoretic-Combinatorial}: Using digit counting and modular arithmetic together
\end{itemize}

\subsection*{C. Domain-Specific Approaches}
\begin{itemize}[leftmargin=*]
    \item \textbf{Number Theory}: Chinese Remainder Theorem, modular arithmetic, divisibility rules
    
    \item \textbf{Algebra}: Not primary approach
    
    \item \textbf{Geometry}: Not applicable
    
    \item \textbf{Combinatorics}: Counting digits, but secondary to number theory
    
    \item \textbf{Probability}: Not applicable
\end{itemize}
\end{tacticalbox}

\begin{techniquesbox}
\subsection*{A. Technique Recognition Index}
\begin{itemize}[leftmargin=*]
    \item \textbf{T18-CRT}: Chinese Remainder Theorem -- PRIMARY technique
    \item \textbf{T19-ModArith}: Modular Arithmetic -- essential tool
    \item \textbf{T20-FLT}: Fermat's Little Theorem -- not needed (45 is composite)
    \item \textbf{Divisibility Rules}: Sum of digits for mod 9, last digit for mod 5
\end{itemize}

\subsection*{B. Technique Selection}
\textbf{Primary Technique: T18-CRT (Chinese Remainder Theorem)}

\textbf{Why it applies}:
\begin{itemize}
    \item $45 = 9 \times 5$ with $\gcd(9,5) = 1$
    \item Problem asks for $N \bmod 45$
    \item CRT allows decomposition into simpler subproblems
\end{itemize}

\textbf{Configuration cues}: Composite modulus with coprime factors

\textbf{Objective cues}: Finding remainder mod composite number

\textbf{Supporting Technique: T19-ModArith (Modular Arithmetic)}

\textbf{Why it applies}:
\begin{itemize}
    \item Need to compute $N \bmod 9$ using digit sum property
    \item Need to compute $N \bmod 5$ using last digit
    \item Modular arithmetic simplifies the large number
\end{itemize}

\subsection*{C. Technique Integration}
\begin{enumerate}
    \item Use divisibility rule: $N \equiv (\text{sum of digits}) \pmod{9}$
    \item Compute digit sum: $S = 1+2+\cdots+44 = 990$
    \item Find $990 \bmod 9 = 0$, so $N \equiv 0 \pmod{9}$
    \item Last digit of $N$ is 4 (from ``44''), so $N \equiv 4 \pmod{5}$
    \item Apply CRT: find $x$ such that $x \equiv 0 \pmod{9}$ and $x \equiv 4 \pmod{5}$
\end{enumerate}

\subsection*{D. Conflict Resolution}
No technique conflicts arise. CRT and modular arithmetic work together seamlessly.
\end{techniquesbox}

\begin{keyinsightbox}
The key insight is recognizing that $45 = 9 \times 5$ immediately suggests using the Chinese Remainder Theorem. We can solve two much simpler problems:
\begin{itemize}
    \item $N \bmod 9$: Use the divisibility rule that a number has the same remainder mod 9 as its digit sum
    \item $N \bmod 5$: Use the last digit of $N$, which is 4
\end{itemize}
Then combine using CRT to get $N \bmod 45 = 9$.
\end{keyinsightbox}

\begin{verificationbox}
\subsection*{A. Verification Level Selection}
This is Problem 19 (late-band), so use \textbf{Level 3: Comprehensive Verification}

\textbf{Decision criteria}:
\begin{itemize}
    \item High difficulty (P19/25)
    \item Multi-step reasoning
    \item Multiple arithmetic computations
    \item CRT application requires care
\end{itemize}

\subsection*{B. Specialized Verification Methods}
\begin{itemize}[leftmargin=*]
    \item \textbf{Computational Checks}:
    \begin{itemize}
        \item Verify digit sum: $1+2+\cdots+44 = \frac{44 \times 45}{2} = 990$
        \item Check $990 = 110 \times 9$, so $990 \equiv 0 \pmod{9}$ \checkmark
        \item Verify last digit is 4 \checkmark
    \end{itemize}
    
    \item \textbf{Boundary/Limiting Cases}:
    \begin{itemize}
        \item Test with smaller concatenation: $N_2 = 12$ has digit sum 3, last digit 2
        \item $12 \bmod 45$: $12 \equiv 3 \pmod{9}$ and $12 \equiv 2 \pmod{5}$
        \item CRT: $x \equiv 3 \pmod{9}$ and $x \equiv 2 \pmod{5}$ gives $x = 12$ \checkmark
    \end{itemize}
    
    \item \textbf{Answer Choice Verification}:
    \begin{itemize}
        \item Check answer 9: $9 \equiv 0 \pmod{9}$ \checkmark and $9 \equiv 4 \pmod{5}$ \checkmark
        \item Eliminate others: 1, 4, 18, 44 don't satisfy both conditions
    \end{itemize}
    
    \item \textbf{CRT Application Check}:
    \begin{itemize}
        \item Need $x \equiv 0 \pmod{9}$ and $x \equiv 4 \pmod{5}$
        \item Form: $x = 9k$ for some $k$; need $9k \equiv 4 \pmod{5}$
        \item $9 \equiv 4 \pmod{5}$, so $4k \equiv 4 \pmod{5}$, thus $k \equiv 1 \pmod{5}$
        \item Smallest such $x$ is $9 \times 1 = 9$ \checkmark
    \end{itemize}
\end{itemize}

\subsection*{C. Confidence Calibration}
\textbf{Confidence Level}: 95\% (High)

\textbf{Reasoning}:
\begin{itemize}
    \item Digit sum calculation is straightforward
    \item Last digit identification is certain
    \item CRT application is standard
    \item All verification checks pass
    \item Answer matches expected pattern
\end{itemize}

\textbf{Potential concerns}: None significant; computation is reliable
\end{verificationbox}

\begin{warningbox}
\textbf{Common Pitfall \#1}: Forgetting that the digit sum includes the digits from two-digit numbers. For example, 10 contributes $1+0=1$, not 10, to the digit sum.

\textbf{Common Pitfall \#2}: Incorrectly applying CRT. The systematic approach is:
\begin{itemize}
    \item From $x \equiv 0 \pmod{9}$, write $x = 9k$
    \item Substitute into $x \equiv 4 \pmod{5}$: we get $9k \equiv 4 \pmod{5}$
    \item Since $9 \equiv 4 \pmod{5}$, this becomes $4k \equiv 4 \pmod{5}$
    \item Therefore $k \equiv 1 \pmod{5}$, giving $k = 1$ and $x = 9$
\end{itemize}

Alternative CRT formula: $x \equiv 9(-1)(4) + 5(1)(0) \equiv -36 \equiv 9 \pmod{45}$, but the substitution method above is more intuitive.

\textbf{Common Pitfall \#3}: Miscounting the digit sum. Double-check: $S = 1+2+\cdots+44 = 990$.
\end{warningbox}

\begin{tipbox}
\textbf{Pro Tip \#1}: Whenever you see a composite modulus, immediately check if the factors are coprime. If so, CRT is your friend!

\textbf{Pro Tip \#2}: For mod 9, always use the digit sum shortcut. For mod 5, always use the last digit.

\textbf{Pro Tip \#3}: In CRT problems, work backwards from the answer choices to verify your result quickly.

\textbf{Pro Tip \#4}: To understand why $N$ has 79 digits: numbers 1-9 contribute 9 digits, numbers 10-44 contribute $35 \times 2 = 70$ digits, for a total of $9 + 70 = 79$ digits.
\end{tipbox}

\begin{solutionbox}
\subsection*{Solution Roadmap}

\textbf{Step 1: Recognize the Structure}
\begin{itemize}
    \item Identify that $45 = 9 \times 5$ with $\gcd(9,5) = 1$
    \item Plan to use Chinese Remainder Theorem
\end{itemize}

\textbf{Step 2: Compute $N \bmod 9$}
\begin{itemize}
    \item Use divisibility rule: $N \equiv (\text{digit sum}) \pmod{9}$
    \item \textit{Why this works:} Since $10 \equiv 1 \pmod{9}$, any power $10^k \equiv 1 \pmod{9}$. Thus a number $d_n \cdot 10^n + \cdots + d_1 \cdot 10 + d_0 \equiv d_n + \cdots + d_1 + d_0 \pmod{9}$
    \item Digit sum $S = 1+2+\cdots+44 = \frac{44 \times 45}{2} = 990$
    \item Note: We count each digit individually, so "10" contributes $1+0=1$, "42" contributes $4+2=6$, etc.
    \item $990 = 110 \times 9 \equiv 0 \pmod{9}$
    \item Therefore, $N \equiv 0 \pmod{9}$
\end{itemize}

\textbf{Step 3: Compute $N \bmod 5$}
\begin{itemize}
    \item The last digit of $N$ is 4 (from the number 44)
    \item Therefore, $N \equiv 4 \pmod{5}$
\end{itemize}

\textbf{Step 4: Apply Chinese Remainder Theorem}
\begin{itemize}
    \item Need to find $x$ such that:
    \begin{align*}
        x &\equiv 0 \pmod{9}\\
        x &\equiv 4 \pmod{5}
    \end{align*}
    \item From the first condition, $x = 9k$ for some integer $k$
    \item Substitute into the second: $9k \equiv 4 \pmod{5}$
    \item Since $9 \equiv 4 \pmod{5}$, we have $4k \equiv 4 \pmod{5}$
    \item Thus $k \equiv 1 \pmod{5}$, so $k = 1$ (taking smallest positive value)
    \item Therefore $x = 9 \times 1 = 9$
\end{itemize}

\textbf{Step 5: Verify}
\begin{itemize}
    \item Check: $9 \equiv 0 \pmod{9}$ \checkmark
    \item Check: $9 \equiv 4 \pmod{5}$ \checkmark
    \item Answer: $\boxed{\textbf{(C)}\ 9}$
\end{itemize}

\subsection*{Alternative Approaches}

\textbf{Method 2 (Subtraction Pattern):}
Another elegant solution observes that subtracting 9 from $N$ gives a number ending in 35 (i.e., $\cdots 4335$), which is divisible by both 9 (digit sum becomes $990 - 9 = 981 = 109 \times 9$) and 5 (ends in 5). Since we subtract 9 to reach a multiple of 45, the remainder is 9.

\textbf{Method 3 (Powers of 10):}
Express $N$ as a sum: $N = 44 + 43 \cdot 10^2 + 42 \cdot 10^4 + \cdots + 1 \cdot 10^{78}$

Key observation: $10^k \equiv 10 \pmod{45}$ for all $k \geq 1$ (since $10^2 = 100 \equiv 10 \pmod{45}$)

Therefore:
\begin{align*}
N &\equiv 44 + 43(10) + 42(10) + \cdots + 1(10) \pmod{45}\\
&\equiv 44 + 10(43 + 42 + \cdots + 1) \pmod{45}\\
&\equiv 44 + 10 \cdot \frac{43 \cdot 44}{2} \pmod{45}\\
&\equiv 44 + 10(946) \pmod{45}\\
&\equiv 44 + 9460 \pmod{45}
\end{align*}

Since $9460 = 210 \times 45 + 10 \equiv 10 \pmod{45}$:
$N \equiv 44 + 10 = 54 \equiv 9 \pmod{45}$

\subsection*{Comprehensive Pitfall Guide}
\begin{enumerate}
    \item \textbf{Digit Sum Error}: Ensure you count all digits correctly, including those from two-digit numbers
    \item \textbf{CRT Misapplication}: Remember to check $\gcd(m_1, m_2) = 1$ before applying CRT
    \item \textbf{Modular Arithmetic}: Don't confuse $N \equiv 4 \pmod{5}$ with $N = 4$
    \item \textbf{Last Digit}: The last digit of the concatenation is 4, not 44
\end{enumerate}
\end{solutionbox}

\begin{competitionbox}
\subsection*{A. Time Management}
\begin{itemize}[leftmargin=*]
    \item \textbf{Estimated Time}: 3-4 minutes for late-band problem
    \item \textbf{Breakdown}:
    \begin{itemize}
        \item Problem reading and CRT recognition: 30 seconds
        \item Computing digit sum: 60 seconds
        \item Applying CRT: 90 seconds
        \item Verification: 30 seconds
    \end{itemize}
    \item \textbf{If stuck}: Move to next problem after 4 minutes; return if time remains
\end{itemize}

\subsection*{B. Problem Prioritization}
\begin{itemize}[leftmargin=*]
    \item This is a number theory problem -- prioritize if you're strong in NT
    \item If you recognize CRT immediately, this becomes a mid-difficulty problem
    \item Skip if you're uncomfortable with modular arithmetic; return later
\end{itemize}

\subsection*{C. Risk Assessment}
\begin{itemize}[leftmargin=*]
    \item \textbf{Risk Level}: Medium
    \item \textbf{Confidence needed}: Must be confident in digit sum calculation and CRT
    \item \textbf{Abandonment criterion}: If you don't recognize CRT pattern within 1 minute, move on
\end{itemize}

\subsection*{D. Strategic Notes}
\begin{itemize}[leftmargin=*]
    \item This problem rewards recognition of standard techniques (CRT, modular arithmetic)
    \item The large number is intentionally intimidating; don't let it scare you
    \item Practice CRT problems to build pattern recognition for composite moduli
\end{itemize}
\end{competitionbox}

\section*{Summary}

This problem exemplifies late-band AMC12 number theory: it appears complex due to the 79-digit number, but reduces to a straightforward application of the Chinese Remainder Theorem once you recognize that $45 = 9 \times 5$. The key insights are:
\begin{enumerate}
    \item Recognize composite modulus $\Rightarrow$ use CRT
    \item For mod 9: use digit sum (standard divisibility rule)
    \item For mod 5: use last digit
    \item Apply CRT to combine results
\end{enumerate}

With practice, problems like this become highly accessible and should take 3-4 minutes to solve confidently.

\end{document}